\documentclass[theme=blue, thai]{cncs}
\RequirePackage{cncs-computing}

\subject{การโปรแกรม 2}
\semester{2}
\academicyear{2568}
\school{โรงเรียนชลกันยานุกูล}

\begin{document}

\worksheet{ใบกิจกรรม}
\studentinfo

\activitysection{ตอนที่ 1: การใช้วงวน}
ให้นักเรียนคาดเดาผลลัพธ์ของโปรแกรมนี้

\begin{cncscode}[C]
#include <stdio.h>
int main(void) {
    int i=0;
    for(i=0; i<10; i++){
        printf("%d", i);
        printf("\n");
    }
    return 0;
}
\end{cncscode}

\noindent เติมผลลัพธ์ลงในตารางตัวแปร (ให้ทำทุก ๆ รอบที่มีการเปลี่ยนแปลง)
\begin{vartable}
     & \\ \hline
     & \\ \hline
     & \\ \hline
     & \\ 
\end{vartable}

\noindent เขียนผลลัพธ์การทำงานของโปรแกรม
\answerlines{4}


\activitysection{ตอนที่ 2: การเขียนโปรแกรมเพื่อแก้ปัญหาดอกเบี้ยทบต้น}
โปรแกรมที่ออกแบบมาจะต้องรับค่าต่อไปนี้\\[0.5em]
$P$ แทน เงินต้น\\
$i$ แทน อัตราดอกเบี้ยต่องวด\\
$n$ แทน จำนวนงวดทั้งหมด

\vspace{1em}
\noindent \textbf{ตัวอย่างเสริม:} หากต้องการให้นักเรียนเขียน Python สามารถใช้โค้ดดังนี้ได้เลยครับ
\begin{cncscode}[Python]
def calculate_compound(p, i, n):
    # เขียนสูตรคำนวณดอกเบี้ยทบต้นที่นี่
    
    return
\end{cncscode}

\drawingbox[ออกแบบ Flowchart ของโปรแกรม]{6cm}

\vspace{1em}
\noindent \textbf{ทดสอบ:} โค้ดที่ไม่แสดงเลขแถว (cncscode*)
\begin{cncscode*}[Python]
print("no line numbers here")
\end{cncscode*}

\noindent เติมโค้ดที่ขาดหายไปในกล่องด้านล่าง
\codeblank[เติม for loop ที่ขาดไป]{3cm}

\vspace{1em}
\noindent \textbf{ทดสอบ:} ช่องเติมคำแบบแทรกในโค้ด (inline blank)
\begin{cncscode}[Python]
import math
(*\blank*) = int(input("Input a : "))
b = int(input("(*\blank[3cm]*)"))
c = math.sqrt(a**2+(*\blank[1cm]*))
print((*\blank[2cm]*))
\end{cncscode}

\end{document}
